\documentclass[letter, 11pt]{article}
\usepackage[utf8]{inputenc}
\usepackage[english]{babel}
\usepackage{amsfonts, amsmath, amsthm, amssymb}
\usepackage[dvips]{graphicx}
\usepackage{url}
\usepackage[top=3cm,bottom=3cm,left=3.5cm,right=3.5cm,footskip=1.5cm,headheight=1.5cm,headsep=.5cm,textheight=3cm]{geometry}
\usepackage{listings}
\usepackage{color}
\usepackage{fancyvrb}
\usepackage{fancyhdr}
\usepackage{enumerate}
\usepackage{multirow}
\usepackage{hyperref}
%% to use the warpfigure capability
\usepackage{wrapfig}
%% Double spacing
\usepackage{setspace}
%% To use appendix
\usepackage[title,titletoc,toc]{appendix}
% Justify caption text
\usepackage[justification=centering]{caption}

%%%%%%%%%%%%%%%%%%%%%%
%Estilo del documento%
%%%%%%%%%%%%%%%%%%%%%%
\pagestyle{fancyplain}

%%%%%%%%%%%%%%%%%%%%%%%%%%%%%%%%%%%%%%%%%%%
%Fancyheadings. Top y Bottom del documento%
%%%%%%%%%%%%%%%%%%%%%%%%%%%%%%%%%%%%%%%%%%%
% Recuerde que en este documento la portada del documento no posee
% numeracion, pero de igual manera llamaremos a esa primera pagina la numero
% 1, y la que viene la dos. Esto es para tener una idea de las que
% llamaremos pares e impares	
\lhead{} %Parte superior izquierda
\rhead{\bf \it UTFSM, Casa Central} %Parte superior derecha
\lfoot{JSCO} %Parte inferior izquierda.
\cfoot{} %Parte inferior central
\rfoot{\bf \it Page \thepage} %Parte inferior derecha
\renewcommand{\footrulewidth}{0.4pt} %Linea de separacion inferior
\renewcommand{\thefootnote}{\roman{footnote}}

\theoremstyle{plain}
\newtheorem{thm}{Theorem}[section]
\newtheorem*{theoremnon}{Theorem}

\theoremstyle{definition}
\newtheorem{defn}[thm]{Definition} % definition numbers are dependent on theorem numbers
\newtheorem{exmp}[thm]{Example} % same for example numbers

\setcounter{secnumdepth}{5} %activate subsubsubsection = paragraph

\begin{document}
%%%%%%%%%%%%%%%%%%%%%%%%%%%
% Tema de las referencias %
\bibliographystyle{acm}
%%%%%%%%%%%%%%%%%%%%%%%%%%

%%%%%%%%%%%%%%%%%%%%%%%%%%
%Definicion de la portada%
%%%%%%%%%%%%%%%%%%%%%%%%%%
\begin{titlepage}
    \begin{center}
	\begin{tabular}{ccc}
	    \includegraphics[width=3cm]{img/utfsm}
	    & 
	    \hspace{-0.2cm}
	    \begin{tabular}{c}
		Universidad T\'ecnica Federico Santa Mar\'ia \\ \hline
		\vspace{0.2cm}
		Departamento de Inform\'atica\\
		\vspace{1.2cm}
	    \end{tabular}
	    \hspace{0.2cm}
	    &
            \includegraphics[width=2cm]{img/di}
	\end{tabular}

	\vspace{1cm}
	%Titulo del Documento
	\begin{tabular}{c}
		\Huge{\sc{Testing, Detection and Possible}}\\\\
		\Huge{\sc{Solutions for the Bufferbloat}}\\\\
		\Huge{\sc{Phenomenon on Local Networks.}}\\\\
		%\Large{\sc{Gemini Observatory}}\\\\
	\end{tabular}

    	\vspace{2cm}
	\Large{\sc{JUAN S. CATALAN OLMOS}}\\
	\vspace{2cm}
	\large{Definici\'on de Tema de Memoria}\\
	\large{para optar al T\'itulo de:}\\
	\large{INGENIERO CIVIL INFORMATICO}\\
	\vspace{2cm}
         	\normalsize{Profesor Gu\'ia: Horst H. von Brand}\\
	%Fecha
    
    \vspace{2cm}
    \normalsize{\sc{\today}}\\
    \normalsize{\sc{Valpara\'iso, Chile}}\\
    \end{center}
\end{titlepage}


\newpage

\tableofcontents
\listoffigures
\listoftables

\newpage

\doublespacing
%section Intro
\section{Introduction}
%If a little salt makes food taste better, then a lot must make it taste great, right?. What happens if you apply the same statement to a network domain? It keeps been as good as it was? It improves the performance or makes it worse?.\\

\textit{Lets think of a network as a road system where everyone drives at the maximum speed. When the road gets full, there are only two choices: crash into other cars, or get off the road and wait until things get better. The former isn't as disastrous on a network as it would be in real life: losing packets in the middle of a communication session isn't a big deal. But making a packet wait for a short time is usually better than ``dropping" it and having to wait for a re-transmission.}\cite{bufferarms}\\

At this point, the role of the router becomes important. It has to control the congestion effectively in networks. It is important to remember that the traffic in a network is inherently bursty, so the role of the buffers in the router is to smooth the flow of traffic. Without any buffering, to allocate the bandwidth evenly would be impossible. But there are some problems with current algorithms; they use tail-drop based queue management that has two big drawbacks: 1.- lockout 2.- full queue that impact with a high queue delay.\\

These problems are fixed with the creation of a group of FIFO based queue management mechanisms to support end-to-end congestion control in the internet. That procedure is called Active Queue Management (AQM). With AQM the loss of package and the average queue length is reduced; this impacts in a decreasing end-to-end delay by drooping packages before buffer comes full, using the exponential weighted average queue length as a congestion indicator. For the proper use of AQM, it has to be widely enabled and consistently configured the router.\\

\textit{Today's networks are suffering from unnecessary latency and poor system performance. The culprit is Bufferbloat, the existence of excessively large and frequently full buffers inside the network. Large buffers have been inserted all over the Internet without sufficient thought or testing. They damage or defeat the fundamental congestion-avoidance algorithms of the Internet's most common transport protocol. Long delays from bufferbloat are frequently attributed incorrectly to network congestion, and this misinterpretation of the problem leads to the wrong solutions being proposed.}\cite{GettysNichols}\\

The existence of cheap memory and a misguided desire to avoid packet loss has led to larger and larger buffers being deployed in the hosts, routers, and switches that make up the Internet. It turns out that this is a recipe for bufferbloat. Evidence of bufferbloat has been accumulating over the past decade, but its existence has not yet become a widespread cause for concern.\\




delay-based and loss-based approach

\cite{NicholsJacobsonCQD}
Cheap memory and ``more is better'' mentality have led to the inflation and proliferation of buffers. Second, dinamically varying path characteristics are much more common today and are the norm at the consumer Internet edge.

as each packet is squeezed down in bandwidth, it must stretch out in time since its size stays constant.

This standing queue, resulting from mismatch between the window and pipe size, is the essence of bufferbloat. It creates large delays but no improvement in throughput. It is not a phenomenon treated by queueing or traffic theory, which unfortunately, results in it being almost universally misclassified as congestion.

The bufferbloat problem, making the window match the pipe size, is hard to address. Window sizes are chosen by senders while queues manifest at bottleneck gateways.
%end intro

\newpage

%section TCP
%subsection{The TCP Protocol \cite{rfc793}\cite{rfc2001}}
%subsubsection{Slow-Strart Algorithm\cite{rfc879}\cite{rfc2460}}
%subsubsection{The Congestion Avoidance Algorithm \cite{rfc2309}\cite{rfc2581}}
%
\section{The Transport Control Protocol}
%http://www.linktionary.com/c/congestion.html
%http://condor.depaul.edu/jkristof/technotes/congestion.pdf 

The Transport Control Protocol or TCP is main transport protocol of the
internet, providing reliable host to host communication over unreliable
transport media\cite{RFC793}. The advantages of its connectionless design,
flexibility and robustness provides reliable, ordered delivery of a stream of
bytes from a program on one computer to another program on another computer,
but for that, the cost are: the needed of careful design so it provides a good
service under heavy loads. Many modifications has been done since it was
originally defined in 1981.

The root idea of any algorithm related to transport connections must be based
into the ``\textit{packet conservation principle}". This principle claims:

\begin{defn}
A new packet isn't put into the network until an old packet leaves.
\end{defn}


\subsection{Fundamentals of TCP}
TCP operates basically making two hosts exchange segments of data. The
connection between these two hosts is identified uniquely by the network addresses and a 16 bit port number at both hosts. The communication between the
hosts is initiated by a three way handshake between them, where the
sequence number is synchronized between the participants. The sequence number
is a 32-bit number and it is the basis for reliable data transport through an unreliable network. This is, starting from the initial sequence
number, each data byte sent as part of the connection has a corresponding
sequence number; and only after having being acknowledged by the receiver is
the data considered to be transmitted successfully.

TCP makes use of the idea of pipe size and the assumption there was reasonable buffering along the data path to send a window of packets at
a time. To control the amount of data that flows through the network path, the
receiver sent the information of how much data it can receive at once, so the
network resource is used efficiently. Window size represents how much data a
device can handle from its peer at one time before it is passed to the
application process. All excess data will just be dropped. This window
is also a constraint to the sender, the sender is not allowed to
transmit more data than the window before acknowledgment of the data sent.

To manage the data sent and the acknowledgment received for those packets, TCP
uses a cumulative scheme. After a packet is in flight from the server with its
corresponding sequence number, sent data goes to a retransmission queue where
it is held until the corresponding acknowledgement from the other end has come in,
or to be resent if not acknowledged within a timeout. When the
acknowledgment of a sequence number is received, the sender discards all data
with sequence numbers bellow the sequence number in the acknowledgment that
has arrived. For retransmission, TCP uses an adaptive scheme. The timeout is
automatically set from the measured round trip time of the connection,
taking into account the variance of the measured values\cite{JacobsonCAC}.
This helps avoid retransmitting potentially lost segments too quickly or too
slowly.

Because today's networks are dynamic and in different configurations, both
topologically speaking as a  bandwidth, TCP must handle these changes and
still be able to maintain communication between the two hosts. The lost of
communications can be an issue that TCP needs to take into account. In case of
lost of one packet means subsequent packets cannot be acknowledged until the
lost packet is retransmitted. This can lad to excessive retransmission and
unnecessary load. TCP extension has been developed that allows the receiver to
send selective acknowledgments of block of received data with sequence numbers
that are not cumulative with the data acknowledged in the traditional
way\cite{RFC2018}.

Since networks are shared and conditions change along the path, the algorithms
continually probe the network and adapt the number of packets in flight. It is
not hard also (but it is often the case) to find along the way decreases in
the bandwidth. This spots are the bottlenecks\footnote{Or the \gls{BW}
points} and they are important because the performance or capacity of the
entire connection (connection as a state between the two hosts) is limited by
the resources that this trace has. With this, controlling the optimal rate of
data transmission is a hard work, and the receiver window as communicated by
TCP is not a necessarily a very good indicator. To fix this issue, TCP
received an addition to its specification: \textit{congestion window}. The
congestion window plays a crucial role in estimating the available bandwidth
between the hosts. After this modification, the minimum between the receiver
window and the congestion window is used as the transmit limit. All the TCP
additions attempt to keep the network operating near the inflection point
where throughput is maximized, delay is minimized, and little loss occurs.


\subsection{TCP's Phases }
To provide a good service, it's needed that TCP flows respond to orders given
by the hosts machines that controls the connection during congestion. This
characteristic of flows is called ``\textit{responsiveness}'' and tells when 
flows must back-off during congestion. In the second half of the 80's, TCP
began to experience a strange phenomenon that manifested itself through a
dramatically diminish throughput. After a deep analysis, Van Jacobson showed
that TCP needed a mechanism to limit transmission speed in the face of
congestion \cite{JacobsonCAC}. This lead to the development of a better
congestion control mechanism for TCP, which was added as a requirement for
hosts connected to Internet \cite{RFC1122}. Since then, the congestion controls
methods implemented in TCP has been updated several times, adding two new
algorithms, the \textit{slow start} and \textit{congestion avoidance} were
designed to keep in ``equilibrium'' the data that is in exchange. Later a
third where added, the \textit{fast retransmit and fast recovery}.

\subsubsection{Slow start and Congestion avoidance}

 \begin{wrapfigure}{r}{0.5\textwidth}
  \begin{center}
    \includegraphics[width=0.48\textwidth]{img/slowstart}
  \end{center}
  \caption[The Chronology of a Slow-Start.]{The Chronology of a Slow-Start. \cite{JacobsonCAC}}
  \label{fig:slowstart}
\end{wrapfigure}

The idea behind bandwidth occupancy in TCP is to use as much as it is allow to
use. With this in mind, slow start proves the network by increasing
exponentially the rate of how many packets sends; the \gls{cwin} is
increases by one packet and adds another for each packet acknowledged. This
behavior continues until the \gls{ssthrsh} is reached or
congestion is signaled. If congestion is detected, the slow start threshold is
reset to half of the amount of congestion window, at the time that congestion
window is reset to a size of one segment, and slow start is restarted. In case
that the slow start threshold is reached, the sender switches to congestion
avoidance mode, which is employed to maintain the transmission rate,
increasing by one segment each RTT\footnote{Each ack increases the congestion
window by $ss*ss/cwnd$. This results in a linear increase of the congestion
window; that means in 1 RTT, congestion window will be increased by
1 \cite{JacobsonCAC} \cite{rfc879} \cite{rfc2460}}.

TCP must be smart enough to signal to the hosts when path is experiencing
congestion, so it must have policies that decrease the network utilization or
increase it if congestion signals are received or not. Also, based on the
afore mentioned principle (``\textit{packet conservation principle}''),
congestion collapse would become the exception rather than a rule. Thus
congestion control involves finding place that violate conservation and fixing
them. It is important to detect early this misbehavior, since new packets are
added exponentially also the congestion grows exponentially, so if is detected
early, only small adjustments will cure it.

Typically, when the sender is signaled that packets are been dropped, TCP/IP
networks understand that as congestion, but typical effects include queuing
delay, packet loss or the blocking of new connections. New implementations
define a explicit congestion notification methods \cite{rfc3168}, that allows
notification of congestion between the end points without any packet loss. TCP
specification mandates initially setting the congestion window to between 2
and 4 segments of data depending on the segment size \cite{rfc3390}.

\subsubsection{Congestion control algorithms}
One of the biggest problems that TCP had in the past years for its development
has been the achieve of optimal utilization preventing congestion associated
with the differences of bandwidth existing in the medium through which
communication is propagated. That is why lately there have been various
implementations for proper handling of the original algorithm 
deployed in TCP for attaining prevent or lessen the effects of this problem.
Among the algorithms that currently exist, the two predominant in modern
operating systems will be mentioned: CUBIC and Compound TCP.

\paragraph{CUBIC TCP}  is the default implementation for congestion control in
the Linux kernel since 2.6.19. \textit{Scaling the window growth rate to match
large \gls{BDP} is rather straightforward, tackling the fairness issues of new
protocols has remained as a major challenge.  Although BDP implies the network
capacity by packet count(or window size), packet count is not adequate to
characterize TCP performance because growth rate of TCP depends on
RTT} \cite{HaCubic}.

Because the basis of their growth is a cubic function (hence its name), has a
less aggressive behavior than their predecessors, while retaining excellent
scalability, fairness and stability. As can be seen in Figure
\ref{fig:cubicfunc}, the cubic function is compound of three parts: a concave
part where the window grows rapidly over small time values, a plateau set at the previous maximal
congestion window size, and the end portion where the increasing bandwidth
above the maximum which exhibits a convex shape.

As can be better seen in the formula at Figure \ref{fig:cubicform}, the congestion
window ain't dependent of its previous acknowledged packets; instead, the
congestion window is computed at each step from a calibrated function. With
this calibration, the plateau is at the previous maximal congestion window
size. This feature allow the algorithm to respond faster to changes in
available bandwidth. Since in previous algorithms throughput is defined by the
packet loss rate as well as RTT, the throughput in CUBIC is defined by only
the packet loss rate.

\begin{wrapfigure}[6]{r}{0.5\textwidth}
    %\rule{5.5cm}{7.1cm}
\begin{center}
    \includegraphics[width=0.48\textwidth]{img/cubic}
  \end{center}
\caption{The window growth function of CUBIC}
\label{fig:cubicfunc}
\end{wrapfigure}

Another feature that CUBIC has with respect to its predecessor in the Linux
kernel, it also changes the behavior of the Slow Start algorithm, replacing it
with

\begin{figure}[ht]
\begin{minipage}{6cm}
\centering
\[ W_{cubic} = C(t - K)^3 + W_{max} \]
    \caption{CUBIC's Congestion Window.\protect\footnotemark}
    \label{fig:cubicform}
\end{minipage}%
\end{figure}%

\footnotetext{C is a scaling factor, \textit{t} is the elapsed time from the
last window reduction. $W_{max}$ is the window size before the last window
reduction. $K=sqrt[3]{W_{max}beta/C}$. And $beta$ is a constant
multiplication decrease factor applied for window reduction at the time of
loss event}

\noindent  one called Hybrid Slow Start (or HyStart) \cite{HaElephants}. 
The idea behind this algorithm is to prevent the extra drop of packets caused by the
overshoot of bandwidth done by the regular slow start. This extra amount of data 
cause unnecessary delay on both ends hosts and network, and taking long time to
recover from, that can be worst if the packet drop is made at a late stage of
the slow start\footnote{since slow start adds one extra packet for each
acknowledged}.

\paragraph{Compound TCP (CTCP)} is the modification to TCP stack algorithm promoted
for Microsoft and implemented in their systems since Windows Vista and Server
2008. This algorithm aims to achieve a significant improvement in resource
utilization by TCP, presenting an fair mix between delay-based approach and
loss-based approach.

CTCP is based in the idea that if the link is under-utilized, the protocol
should increase the sending rate in a aggressively manner so it can obtain
available bandwidth more quickly. But when the link is nearly full
utilization, keep that aggressiveness is unproductive because at this stage it
is expected that the protocol is more accurate in achieving complete the
remaining bandwidth. So CTCP incorporates a scalable delay-based component
into the standard TCP congestion avoidance algorithm \cite{4146841}.

\begin{wrapfigure}{r}{0.5\textwidth}
    %\rule{5.5cm}{7.1cm}
  \begin{center}
    \[ win = min(cwnd + dwnd, awnd)\]
  \end{center}
  \caption{Compound TCP sending window.\protect\footnotemark}
  \label{fig:ctcpform}
\end{wrapfigure}

As can be seen in figure \ref{fig:ctcpform}, the congestion window or cwnd
which manages the loss-based component is untouched, and a delay window - dwnd
- is added to manage the delay-based component.  The awnd is the advertised
window from the receiver. This algorithm has shown that in periods of under-
utilization, it has an aggressive and scalable increase rule, and gradually it
reduce the sending rate accordingly when the network is sensed to be fully
utilized. But thanks to the packet loss component, it also reacts to packet
loss reducing the sending rate to keep the conservative property of TCP. Also,
Tan et al \cite{Tan06compoundtcp} has reveal that CTCP has a good performance
over current high BDP networks too.

\subsubsection{Fast retransmit and fast recovery} 
Timers are an important part of this self-regulated protocol. TCP maintains
four timers\footnote{The timers that TCP has for each connection are:
(1)Persist Timer: Ensures that window size information is transmitted even  if
no data is transmitted. (2)Keepalive Timer: Detects crashes on the other end
of the connection. (3)2MSL Timer: Measures the time that a connection has been
in the TIME\_WAIT state. (4)Retransmission Timer: The timer is started during
a transmission. A timeout causes a retransmission}. One of them, the
retransmission timer, causes the sender to retransmit the segment over which
at the expiry of this timer(which is a function of the estimated RTT), has not
yet received their respective acknowledge, assuming the segment is lost in the
network. When this action is unleashed, the server will interpret it as a sign
of congestion, but data still flowing.

As already seen in \textbf{2.2.2.1}, slow start can do a lot of damage to our
networks, but fast retransmit helps to recover from lost without hurting the
network that much. It is triggered after a specified number of
acknowledgments, usually is at the third, with the same acknowledge number are
received by the sender. This makes the sender reasonably confident that
the next higher sequence number was dropped and needs to be retransmitted; all
this without the need for the action triggered by the timeout.

Fast recovery manages the continuous transmitting of new data on subsequent
duplicate acknowledgments, until the retransmitted packet is received and
after a non-duplicated acknowledge is received, the congestion window is reset
back to the value it had before the fast retransmit was initiated, and the
fast recovery mode is ended, going back to congestion avoidance mode.

%end TCP

\newpage

%section Buffers 
\section{Router's Buffers effects}
The main function of routers is to absorb bursts of traffic coming from the
routers, and ensure that links are used to their maximum capacity. A network
that does not have buffers, do not have the ability to store packets that are
waiting to be transmitted, so any package that is over the capacity of the
link will be dropped. In order to operate without buffers, arrivals must be
constant and predictable; synchronization of the entire network would be
necessary to avoid any loss. Such networks are complex in design, expensive to
implement and particularly restrictive.

The throughput of each network is limited by the capacity of the slowest link
causing a bottleneck; no matter how much more packets are introduced to the
network, the transfer will not be faster than that determined by this link.
That is why the major location of buffers in areas where bottlenecks occur is
essential for proper functioning of the network, but this transfer between
fast-to-slow networks is different on each route, even for the reverse path.

The presence of buffers is necessary to help to reduce data loss. Also because
in the past the high cost of memory kept the buffers quite small, causing it
to fill pretty quickly after the network became saturated, signaling to the
communications protocol the presence of congestion and thus the need for
compensating adjustments.

Unfortunately the network transport protocols, in order to operate at full
capacity, require that the hosts are notified in a timely manner when they
should back up and thus able to adapt its transmission rate to the available
capacity. The absence of such timely notification triggers the presence of
full buffers and increased communication latency.
\subsection{Full Buffers}
\subsection{High Latency}
\subsection{Sizing Router's Buffers}

%end Buffers

\newpage

%section AQM
\section{Active Queue Management}
Ensure communication only checking the endpoints has revealed that it is not
enough and that it is also necessary to check the status of in-flight data
transmission. With the passage of time and with increasingly high speed
networks it has become more and more important to have mechanisms that keep
throughput high but average queue size low. The changes implemented in TCP for
congestion control have proven not to maximize the capacity of the medium
along with a performance degradation. Among the most common problems are:
connections still experience multiple packet loss, low link utilization and
congestion collapse. It is important to keep in mind that when a packet is
dropped before it reaches its destination, all of the resources it has
consumed in transit are wasted.

This is how the role of routers extends beyond joining and interconnecting local
networks and the Internet together. The role of routers becomes important
because they are able to better manage congestion that may occur in networks
due the knowledge they have about other routers in the network and can choose
the most efficient path for the data to follow, or as already seen, send
message to the ends to fall back or even drop data.  They must allocate
the available bandwidth fairly between all flows.

Typically the queues that routers maintain are designed to smooth and
accommodate bursts of data delivered by the hosts and transient congestion,
but as the queues start to fill, the routers drop packets using a FIFO based
drop-tail management. This discipline has two mayor drawbacks:
\begin{enumerate}
\item Lock Out: a small number of flows tends to monopolize the usage of the
buffer capacity\cite{evolvshortlongflows} .
\item Full Queue: these buffers tends to be always full, leading to 
high latency.
\end{enumerate}

Also, as explained in \cite{FloydJacobsonRED}, when a router use tail-drop,
\emph{the more bursty the traffic from a particular connection, the more
likely it is that the gateway queue will overflow when packets from that
connection arrive at the gateway/router}.

All this led to develop and implement more aggressive or active strategies to
congestion control on the Internet called Active Queue Management or AQM for
their acronym. AQM is a group of FIFO based queue management mechanisms to
support end-to-end congestion control in the Internet. The goals behind AQM
are to reduce the average queue length and with that decreasing the end-to-end
delay. Also to reduce packet losses that reflect as a more efficient resource
allocation.

AQM maintains the network in a region of low delay and high throughput by
dropping packets before queues become full and can reliably distinguish
between propagation delay and persistent queuing delay. Also, because the
router can monitor the size of the queue over time, the gateway is the
appropriate agent to detect incipient congestion. The most common AQM
algorithms are: RED, SRED, BLUE, SFB, CoDel.

\subsection{RED}
Random Early Detection or RED is the algorithm introduced in
\cite{FloydJacobsonRED}. RED monitors the average queue size and when it
exceeds a preset threshold, the router drops the arriving packet with a
certain probability, where the exact probability is a function of the average
queue size. If \gls{ECN} is active, the packets are choosen randomly and marked. When
a connection uses a larger share of available bandwidth, it is more likely to
be marked. The idle periods (periods when the queue is empty) are also taken
into account for the calculations of the average queue size.

The objectives behind the  implementation of RED are:
\begin{enumerate}
\item Minimize packet loss and queue delay.
\item Avoid global synchronization of sources by randomizing its marking
decision and by spacing out its marking.
\item Maintain high link utilization.
\item Remove biases against bursty sources by limiting the queue occupancy so
that there is always room left in the queue to buffer transient.
\end{enumerate}

\begin{wrapfigure}{r}{0.5\textwidth}
  \begin{center}
    \includegraphics[width=0.48\textwidth]{img/RED}
  \end{center}
  \caption{RED's Flowchart}
  \label{fig:RED}
\end{wrapfigure}

Because RED can control the average queue size while accommodating transient
congestion, implementing RED in routers can provide high throughput and low
average delay in high-speed networks with TCP connections that have large
windows. Also, in routers where RED and ECN are implemented the bursty traffic
is well handled and the global synchronization between flows is avoided by
decreasing their window at the same time.

The most common criticism to RED is its complexity to properly configure its
parameters. It is because of this that after their official publication, it
was necessary to perform a second publication \cite{NotesRED}, which aims to
explain a few more details on how to configure it properly. Since RED relies on
queue length as index of congestion, it can reflect the presence of congestion
but not the severity nor the numbers of competing connections sharing the
link. Since long ago networks ceased to ave a static configuration, RED requires
a wide range of parameters to operate correctly under different congestion
scenarios. Unfortunately, this failure leads to deterministically mark packets
and suffer from poor link utilization. \emph{While  ECN timeouts allows RED
to avoid packet loss, the deterministic marking eventually causes
underutilization at the bottleneck link}\cite{FengBLUEAQM}.

As stated in \cite{NicholsJacobsonCQD}, \emph{RED was simple and can be
effective at reducing persistent queues, little guidance was available to set
its configuration parameters and it functioned poorly in a number of cases.
This led to a general reluctance to use it}.

Also to archive the ideal operating point, it can only do so when it has a
sufficient amount of buffer space, which leads to default induce an extra
latency.

\subsection{BLUE}
\begin{wrapfigure}{r}{0.5\textwidth}
    %\rule{5.5cm}{7.1cm}
    \centering
	\begin{verbatim}
	Upon packet loss
	    if (now - last_update >freeze_t)
	        Pm = pm + d1
	        last_update = now
	Upon link idle
	    if (now - last_update >freeze_t)$
	        Pm = pm - d2
	        last_update = now
  	\end{verbatim}
    \caption{BLUE's algorithm}
    \label{fig:BLUEAlg}
\end{wrapfigure}

Introduced by Wu-chang Feng et al in \cite{FengBLUEAQM}, BLUE aims to manage
congestion by using packet loss and link utilization history. The concept
behind BLUE's development is to avoid drawbacks of RED like parameter tuning
problems or to determine the fluctuations of actual queue length. The key idea
is to perform queue management based directly on packet loss and link
utilization rather than on the instantaneous or average queue lengths.

Through the use of a single probabilistic indicator called $p_m$, it
simplifies the process of dropping or marking packets when they are enqueued.
Against a continuous dropping of packets due to overflow, the probability
increases, increasing the rate at which it sends back congestion notification.
Also, to control how this parameter changes over time, the
$freeze\_time$ is used to determine the minimum time interval between
updates.  The value on which to increase or decrease, if the link is idle, is
determined by $\delta_1$ and $\delta_2$.

The effects of BLUE, as proved in \cite{FengBLUEAQM}, are the reduction of the
buffer size, which also reduces the end-to-end delay and helps to improve the
effectiveness of the congestion control algorithm. BLUE has also shown that it
has a better performance when marking incoming packets, leading to congestion
notifications causes periods of sustained packet loss nor periods of continual
underutilization.

\subsection{CoDel}
After several years using algorithms based on the average queue size, Nichols
and Van Jacobson realized that this approach although they have information of
congestion, no hard data on how serious is this. Jacobson et al, determined
that it may submit two types of queues: one that smooths bursts of arriving
packets turning it into a stable and continuous sent and tends to last less
than one RTT (called good queues); and those that only tend to create
excessive delays and lasting several RTTs (called bad queues). With this, they
were able to determine that the core of bufferbloat detection problem is
separating good from bad queues. Bearing in mind the above is that an
algorithm called \textit{Control delay} or CoDel by its acronym was
developed.

Presented in \cite{NicholsJacobsonCQD}, CoDel is based on the idea that it is
sufficient to maintain a threshold above which the tail can not exceed rather
than keeping a window of values to compute the minimum. Rather than measuring
queue size in bytes or packets, it is used the packet-sojourn time through the
queue. The time each packet spends in the queue is independent of the
transmission rate, the \gls{sojourn} is more valuable information about the behavior
of the buffer, in addition to being much more related to the effects the user may
experience. The sojourn is based on a timestamp corresponding to when the
packet arrives at the queue, at which this mark is added to the package
information. The minimum packet sojourn can be decreased only when a packet is
dequeued.

Another difference that occurs is that CoDel accepts the existence of queues,
but it is unacceptable to drop packets when you have a buffer with less than
\gls{MTU}'s bytes.

\begin{wrapfigure}{r}{0.5\textwidth}
    \centering
	$f(n) = \frac{100}{\sqrt{n}}*f(n-1)$ \\
	f(0) = 1
    \caption[CoDel droptime interval]{CoDel droptime interval, n=iteration}
    \label{fig:CoDeldroptime}
\end{wrapfigure}

But when the queue delay has exceeded target for at least interval, a packet
is dropped and a control law sets the next drop time accordingly to Figure
\ref{fig:CoDeldroptime}. The ways in which CoDel stops packet drop is either
when the delay queue is less than the target value or when the buffer contains
less than MTU's bytes. Furthermore the algorithm has additional logic that
controls reenter to dropping state too soon.


Beside the minimum packet sojourn, the other two parameters that CoDel needs
are the target and interval, which are almost self-explanatory. Target is the
acceptable standing queue delay and interval works as the  time on the order
of a worst-case RTT of connections through the bottleneck. The recommended
target value is 5ms since tests reveal that with this value, the utilization
of the links is the optimal. Interval is loosely related to RTT since it is
chosen to give endpoints time to react without being so long that response
time suffers. A setting of 100ms works well. \textit{CoDel's efficient
implementation and lack of configuration are unique features that make it
suitable for managing modern packet buffers}\cite{NicholsJacobsonCQD}.

%end AQM

\newpage

%Flaws
\section{Hidden Flaws and Bufferbloat}
Because much of the basis for how TCP was designed in the early days of the
Internet, where speeds were low, the burden on the networks were only a few
kbps, and the networks interconnecting were separated by a couple
milliseconds, the use was ideal. But it is not hard to see that in the past
few decades, Internet growth has exceeded the expectations of even the most
experienced analysts and dreamers.

Intensive-consumption applications like on-line games, audio and video
streaming for music or on-line calls, torrent or P2P applications and 
sites that have a high level of resource utilization like Youtube, Facebook or
Netflix, have made networks of today quite different from those that formerly
existed, but still running the same protocols\footnote{While many have been
updated, the basics remain the same}. If that was not enough, the creation and
incorporation of new devices such as smart phones, tablets and even smart-watches
with the ability to transmit data over the Internet, makes it more necessary to handle data traffic in an efficient way..

TCP relies on timely notification to adjust its transmission rate to the
available bandwidth, which is commonly signaled by packet loss. But since a
couple of years, and helped by lower prices of hardware along
with the thought that more is always better, the manufacturers have decided to
prevent the loss of data as much as possible with the addition of larger buffers on
their devices. This simple action is slowly bringing a new collapse not only
TCP but data transmission in general, as can be seen in \cite{CACMStaff}

With the increase in the size of buffers, the packages spent more time
``on the fly'' in big buffers instead of dropped, signaling to TCP to
reduce the sending rate. When finally a large amount of data is dropped,
it causes TCP to
sharply drop in transmission rate, freeing up bandwidth. Unfortunately, due to
the size of the buffers being static, when the new TCP slow-start phase starts,
again due to the lack of timely signals, TCP will work with reduced
performance. This problem is cyclical, resulting in exponential TCP back-off,
throughput degradation and very high latencies.

But the paths on Internet are shared by multiple TCP streams so the buffers,
and this back-off behavior, has a tendency to synchronize flows
\cite{main:ref:1}. This causes all the flows to throttle back their
transmission rates simultaneously, amplifying the effect. This decrease in
latency and reduced throughput are the effects of Bufferbloat phenomenon.

As stated in \cite{GettysNichols}, \emph{long delay from Bufferbloat are
frequently misattributed to insufficient bandwidth and this misinterpretation
of the problem leads to the wrong solutions being proposed}. As seen in the
past sections, the packet loss for TCP is not in itself a problem, conversely
it is essential for its functioning. The real problem is  the excessive and
consecutive packet losses that occurs when the buffers kept persistently
fulls. This effects may cause that modern applications timeout their
connections or in others may experience some delay. This is the real
effects and problems of Bufferbloat.

%end Flaws

\newpage

%section Test
\section{Experimental Work}
The goal of the experiments outlined in this section is to examine different 
residential and public networks looking for checking the existence of the 
phenomenon under a uncontrolled scenario. This will be done by running a set of 
tests with different tools first to define the characterize and then measure the 
behaviour of the networks with and without any load. More specifically, latency 
under load is tested and analyzed to identify if the latency that occurs is due 
to an excess of buffers or due to some other problem.\\

This section aims to explain the setup that will be used to perform the tests,
the different tools used in each of these tests and what is to be achieved and 
expected from each one of these tests.\\


\subsection{The Test Setup}
The tests were ran under a pseudo controlled environment, using one
physical  machine and a second virtual machine hosted on the first machine. These
systems runs under a regular OS without any modifications. Also, for some tests two other devices will be
added,  one acting as an Iperf Server and a regular Android Tablet that will
be used to add some extra load to the network when the Ethernet cable is used
as medium.

All of these tests will be carried out in a real-world scenario, where no
packet  prioritization is done by the server against our flows, the routes can
vary between each iteration of the same test, and many different flows will
collide with other flows from different sizes and types. Nor there is more
information about how the flows are treated by the queue manager algorithms or
about how they are  configured.

\subsubsection{Hardware Characterization}

\begin{description}

\item [Physical Machine]: \hfill \\
The physical machine runs as host OS, Windows 7 SP1, that works with an Intel(R)
Core(TM) i7-2670QM CPU \@ 2.20GHz with 8GB of available RAM. This machine will 
always be connected through its wireless adapter, a \textit{Broadcom Corp. 
BCM4313 802.11b/g/n Wireless LAN Controller (rev 01)}. The Ethernet controller 
is a \textit{Realtek Semiconductor Co., Ltd. RTL8111/8168 PCI Express Gigabit 
Ethernet controller (rev 06)} adapter, and will be bridged to the virtual 
machine for some tests.

\item[Virtual Machine]: \hfill \\
The virtual machine is hosted using VMWare Player 6.0.1, with 4 processors 
assigned for use plus 4GB of RAM, and with the Ethernet adapter connected only
for certain tests. The OS selected is a Debian based OS called Kali Linux, and
using the kernel release identified as Debian 3.12.6-2kali.

The wireless adapter is a AIR-802 USB adapter with Zydas chipset and a TP-link
8dbi antenna. This adapter is directly connected to the virtual machine and
hooked to the physical machine without the USB extension, this way any extra
signal  loss is avoided. 

\item[Iperf Server]: \hfill \\ 
This is a VPS hosted by Digital Ocean
\footnote{\url{https://digitalocean.com}} with 512MB Ram, 20GB SSD Disk, and
located in New York data center. This machine runs Ubuntu 12.04.3 x64 under
KVM software using as a processor an Intel Hex-Core 3 GHz.

\end{description}

The idea of using an external device to connect the virtual machine to the 
network and not using a bridged configuration provided by software, is mainly 
because with an USB device the machine will take care of all the management and 
administration of the device, avoiding any possibility that the host machine 
modify or manages any flow. Also, by using a second machine to overload the 
uplink, causes to avoid overflow the queue on the machine that is performing 
the tests. With this, it is expected to minimize the possibility that our 
testing machine is causing extra latency, either by the saturation of the 
wireless channel, or by the queue in the traffic control subsystem into the 
kernel. 

For the tests, the Bufferbloat community\cite{bloat} has created a set of best 
practices\cite{tg12} to follow so the results are consistent and repeatable, 
but in this case, computers and routers will not be modified as it will attempt 
to analyze what an every-day-user experience. Those practices will be taken in 
consideration if it is the QOS present in routers and deactivate as recommended.
 


\subsection{Tools Definition}
The tools used for the benchmark were selected by the capability to determinate 
the presence of the Bufferbloat phenomenon in the network of study by the 
analysis of its indicator, the latency or the round trip time\footnote{To 
Subtel, organization responsible for control and supervision in the performance 
of telecommunications in Chile, RTT and Latency correspond to the same thing 
\cite{subtel} }. So, with this as a main consideration, the tool's selected for 
the benchmark will be chosen by the complexity and accuracy to measure and 
determine the RTT into a IP/TCP connection, and how hard is to consistently 
replicate the results under similar contexts. Other tools were selected to help 
to determine to characterize and proof the capability of the network to cause 
Bufferbloat.\\

The tools selected to be used in this tests are the following:

\begin{itemize}
    \item Speedtest test by Ookla
    \item Netalyzr by ICSI
    \item Iperf Tool with Tcptrace/Xplot.org
    \item Page Benchmarker extension for Google Chrome
    \item Smokeping Latency Tool
\end{itemize}

\subsubsection{Speedtest}

Developed by Ookla\footnote{\url{https://www.ookla.com/about}}, this tool is used
for most of the ISP's and many users in Chile to test their broadband's
connections globally. It can be used not only in their website 
\href{http://www.speedtest.net}{\textit{www.speedtest.net}} but also in Android,
iOS or Windows Phone. A command line interface developed in Python can be used 
too for testing internet bandwidth.\\

The server that will be selected to perform every test, either has or not the
fastest ping, is the one hosted by the Pontificia Universidad Cat\'olica de
Valpara\'iso (this host is the default selected most of times by the site also).  

\subsubsection{Netalyzer\cite{netalyzr}}

asdf

\subsubsection{Iperf}
Iperf is a well known and commonly used network testing tool. It can create a
TCP and UDP data streams and measure the bandwidth and the quality of a network
link. It can perform multiple tests like Latency, Jitter or Datagram Loss.\\

Iperf basically tries to send as much information down a connection as quickly
as possible reporting on the throughput achieved. This tool is especially useful
in determining the volume of data that links between two machines can supply.
This two machines define the network, one acting like a server and the second as
the client. For this scenario, the server will be the VPS that only will receive
the Iperf connections (also will be running ssh but without further
interaction). The VM Linux machine will work as or Iperf Client.\\  

As mentioned in Iperf users mailing list \textit{
When one runs TCP tests, there are 2 things that block Iperf from having clear
view of real throughput: buffering on sender's side (TCP/IP stack) and TCP
behaviour itself (acking). What Iperf can measure is the pace with which it
sends data to TCP/IP stack; TCP/IP stack will only accept data from application
when buffers are not full. If the buffer is huge, Iperf will see high
throughput initially, then it will drop. If there's congestion or retrasmission
going on, Iperf will see it as lower throughput}\cite{iperfmaillist}, but the
data generated by Iperf won't be further analyzed because the  idea behind using
this tool is a TCP's packet generator. This means that the packets generated by 
Iperf will captured and analyzed with tcpdump, tcptrace and xplot.org.\\

\subsubsection{Page Benchmarker}

Page Benchmarker is a Google Chrome extension that intent to test page load 
time performance within Chrome. Measures time-to-first-paint, overall page load 
time KB read/written, and several other metrics, and with its capability to 
clear the cache and existing connections between each page load, makes this 
tool one of the main sources results.\\

\subsubsection{Smokeping}
SmokePing is a latency logging and graphing tool that consists of a running 
daemon which organizes the latency measurements and a CGI which presents the 
graphs. With SmokePing give us the ability to measure latency and packet loss 
in the current network, with RRDtool is capable to maintain a long term 
data store and to draw different graphs with the giving up to the minute 
information on the state of each network connection.\\

Smokeping can be configured to perform a wide range of latency measurement 
probes each one directed to an independent target or over a set of targets 
selected for each proof.\\




\subsection{Test Description}
To test the existence of the Bufferbloat phenomenon, five tests are conducted 
as described below. Each test will be repeated under the following contexts:

\begin{enumerate}
\item Twice on the same day in one network to determine if does exist a 
considerable variance in latency for different times in a day.
\item Select different public and private networks with different 
\textit{``speeds''}.
\item Use the Ethernet cable in order to compare the results with those 
previously obtained using Wireless.
\end{enumerate}

Because it is intended to prove the existence of the phenomenon under
circumstances experienced by everyday users, no kernel parameters will be
amended nor changed, and only the  ability to perform QOS on routers that have
this feature is disabled.  Furthermore, the overall question these tests seek
to answer is the following:

\begin{theoremnon}[] 
``The networks that we use every day, have the necessary characteristics to
generate the  Bufferbloat phenomenon whether under low loads and if does
exists, the how  serious are the effects ?''
\end{theoremnon}

\subsubsection{Speed test} 
The idea under this test, is to set a baseline by comparing the speed offered 
by the ISP and the one at the moment of testing. To find the speed provided, 
the  tool used is the Speedtest in the web site of Ookla.

The benefits of using this web site is that not only it will reveal only
the available national uplink and downlink, also it will set the baseline
for the ping. The expected pings, independently of the connection bandwidth,
should be around $\sim12ms$, based on the data presented in \cite{netvtr} and
\cite{netmov} by the two most used ISP in Chile\footnote{Telef\'onica has a
39.2\% market share while VTR owns 38.8\% respectively. \cite{ispshares}}.

As data, it is estimated that the ratio of the average speed as compared to the
rated speed of the service offered for domestic bonds has a variation to 85\%,
while for international links decreases drastically to 35\%\cite{modcompqos}.
Due this factor, none of the tests will be performed on servers located outside
our country, as it is believed that the data will not very representative
regarding actual service.

\subsubsection{Signs of trouble} 
After characterize our network based on the
speed, it is needed to try to define the state of the service based on the
quality in which our network operates, in example: what kind of traffic is
passing through our network, and also get a little more detail on the
average rates of delay and buffers with which the traffic can find along the
path. To collect all this information will be used Netalyzr. 

While laws in Chile assure consumers networks free of traffic shaping
barriers and filters, ISPs apply traffic management
measures\cite{shapemov}\cite{shapevtr} to deliver an optimum experience of the
service's use and thus, achieve an efficient use of the network, thus further to
protect the safety of users while maintaining network stability.

This is why it is expected that some ports or services are partially or
completely blocked, but the most important is that this test will first light on
the existence of the phenomenon under study.

\subsubsection{Collapse test} 
Already having a clear idea of the state of the studied network, the next will
be try to check empirically that the results obtained by the previous test
above described are valid. For this, Iperf will be used to generate as much
traffic as possible for 5 minutes between servers. But the main goal is not to
measure the throughput of the network; instead to capture the packets that
Iperf produces and study them. So before run Iperf, tcpdump will capture this
information. Subsequently, the tcp's trace is taken and extracted with
tcptrace tool and generate the RTT graph.

This exercise will be conducted three times as part of the test, running the
first time as the sole source of network load. For the second and third time,
after 50 seconds after Iperf was started, from the Windows machine will be 
performed an upload to a Dropbox account with a file big enough to the upload 
take more than the time defined to this test. This upload intends to saturate 
the upstream link, leading to an overload of existing buffer (indistinct 
whether the server is national or not, there are several levels where the 
route is common). The upload will be stopped 50 seconds before the time limit.

The time gap established before loading and the end of the test is to let
Iperf frames reach a steady state flow and thus to generate a basis on which
to analyze the time when the traffic is maximum.

The expected RTT are around $\sim12ms$ and $\sim100ms$ without load, for
national and international servers respectively. About the expected value under
load, it is hard to estimate, further that their expected behavior is 
to be stable round a certain value. In case of no stable behavior, the possible 
causes will be analyzed and tracked.

\subsubsection{Load benchmark test}
With the effects of excess buffers already determined, this tests seeks to show
how this phenomenon affects users who surfs through a web browser. This is why
the extension of Google Chrome Page Benchmarker comes in and it will be used to
do 5 iterations of 10 loads to a website and analyze the behavior of these. 

As in the previous experiment, the first will be without any load on the
network. Then, again a file will be uploaded to Dropbox from the second machine
and after 30 seconds from the start of the load, and with the load active, will
proceed to a second iteration. The third will begin the after a minute,
canceling the file upload and after waiting 30 seconds after cancel the upload,
the benchmark will be measured again. The fourth and fifth will be a minute and
a half of a minute since the previous iteration have finished. The site chosen is
\url{http://www.usm.cl}.

As like Jim Gettys showed in his video demonstration\cite{gettysex}, the 
benchmark increases between 10 to 15 times with load against the original times. 
For this test, the expected proportions are lower, this mainly because the 
kernel machine that will run the tests has implemented already some improvements 
to mitigate this problem.

\subsubsection{Smoke the path}
To verify that the effect of Bufferbloat, if it is the case, does not just
happen on a single server or with a single TCP flow, the Smokeping tool will be
configured to perform two types of tests: Fping and echopinghttp. These tests
will be directed to different types of servers (ie: physical and VPS machines)
which have different connection settings and/or type of services (some servers
has dedicated link), and located in both Chile and the United States.

The tool will be left couple of minutes to track and save the state of the
connection under no extra load (only the test machine will be generating
data). After that, a upload of a big file will be performed from a second
machine until it finish or the behavior has become stable and then re-enter to
a phase without load for couple of more minutes and repeat the upload.

With this test, in addition to verifying the validity of the previously captured
data is expected to determine the presence of other factors that can contribute
to the degradation of the quality of service on the network, whether factors
such as quality of service provided by the host, channel/path problems, problems
related to the way handling packets by the router or the machine, and/or any
other that may arise.

ia	
%end section Test
\newpage
%section Results
\section{Results}
This section aims to show and summarize the results obtained in the different tests detailing each specific case, and then collect and summarize all.


\subsection{Speed test}
As expected, in all cases the networks were asymmetric, 
the download bandwidth (the downlink) being significantly higher than the uploading
or uplink. This also explains the relationship between the uplink and ping,
presenting a major ping in the networks with lower uplink. The relationship
can be better be seen in the Figure \ref{fig:speeds}.

\begin{figure}[ht]
\centering
	%\rule{5.5cm}{7.1cm}
    \includegraphics[width=0.9\textwidth]{img/speed_graph}
\caption{Speed and Pings based on Appendix \ref{app:bwmeasures}, Table \ref{table:speeds}}
\label{fig:speeds}
\end{figure}%

The Figure \ref{fig:speeds} also is a proof of the common missundersanding
about the how the ISP's promote high download speeds but with low uplink. Even
though most of the time home users download content, in order to keep a good
ratio the data needs to flow both ways, due to the way TCP
works, you need good communication in both directions. While in theory this
works without major problems, unfortunately in reality users download content
from different sources at once, and if we add new requirements like real-time  or on-demand applications and online games, we see that not only a
high-speed downlink is important.

Another important point is the growing need for higher bandwidth.
Unfortunately it is common to associate bandwidth with speed. Instead,
bandwidth is a measure of capacity, not how fast the network responds. So,
more capacity is only better if the existing latency is low. This may
explain why the tests performed in network tagged as \emph{``casa''} can have low pings
even when its capacity is close to the mean, but this network is through
fiber\footnote{In practice, fiber service have a propagation delay 10 times
faster than ASDL service}.

\begin{table}[ht]
\begin{center}
\begin{tabular}{|c||c|c||}
 \hline
& \multicolumn{2}{|c|}{Ratio (\%)} \\ \hline
Location	& Uplink		& Downlink \\ \hline \hline
casa\_w		& 99,00			& 103,20 \\ \hline
casa2\_w	& 101,80		&  89,73  \\ \hline
casa\_c		& 103,00		& 102,93  \\ \hline
marcia\_w	& 166,00		&  63,00  \\ \hline
marcia\_c	& 166,00		&  62,90  \\ \hline
sara\_w		& 112,00		&  54,50  \\ \hline
sara\_c		& 114,00		&  80,80 \\ \hline
polmos\_w	& 116,00		&  49,03  \\ \hline
polmos2\_w 	& 117,00		&  46,80  \\ \hline
nalucem\_w	& 114,00		&  99,25  \\ \hline
nalucem2\_w & 112,00		& 101,50  \\ \hline
mbahamon\_w & 120,00		& 102,10  \\ \hline
4low\_w	 	& 120,00		&  80,30  \\ \hline
hidalgo\_w 	& 119,00		&  89,80  \\ \hline
\end{tabular}
\caption[Speed Test: Variation ratio of offered vs measured speed ]{Variation ratio of offered vs measured speed.}
\label{table:ratiospeed}
\end{center}
\end{table}

The $\sim12ms$ ping were also not obtained, but the values ​​are still within
the acceptable range around the $\sim26ms$. These values ​​may be due to
factors such as congestion, conditions on the server, or the processing time
information. The variation of the ratio between the measured and the bandwidth
contracted is of 80\%, which is only 6\% less than the spected
value\footnote{Details of the provider and bandwidths can be found on Appendix
\ref{app:bwmeasures}, Table \ref{table:comparative}}.

Further analysis is needed to determinate how the bandwidth is related with
the Bufferbloat, but it is clear that less the bandwidth, higher the ping, but
it can not be stated the oppossite, higher the bandwidth, less the ping
because it also depends on how fast the network is and the effects on the
latency \cite{main:ref:3}.


\subsection{Netalyzr test}
After having a clear idea about the capacity of the network, it is necessary
to know how it speed and behavior works. This requires to analyzing how
latency behaves. As stated, latency is the time it takes a message to travel
from one computer to a sever, and has a huge impact on how the user experience
the network.

Within all the information that Netalyzr provides, it only will be taken only
the most relevant information to determining whether Bufferbloat is presented
on our network. For this, the data that will be take into consideration is:

\begin{description}

\item [DNS resolver Time:] This test intent to measure how quickly the DNS
resolver is able to resolve the mnemonic name. Slow resolvers may make the
network seem ``slow'' even if the network itself is fast.

\item [Network Buffers:] The most revealing fact of the existence of the
phenomenon studied. While it is a fact that is commonly overlooked, also it is
crucial to the quality of the network. If the buffer is too small, network
protocols such as TCP are unable to send as fast as the network allows. If the
buffer is too large, a single transfer will fill up the buffer, delaying all
other traffic.

\item [Network Performance:] To determine the performance of the network,
Netalyzr measures latency by sending a series of small messages to our server
and then seeing how long it takes the messages to return. Since the
communication if from outside US, the latency should naturally be higher, but
how much higher?
\end{description}


The results in table \ref{table:buffer} are much more clarifying than those
obtained in the previous test. In the first data set related with DNS, the
times are relatively stable and the time that took to resolve the requests
were almost indiscernible. While it is common to find cases with shorter at
$\sim15ms$ for local name resolution\footnote{ As example can be a site hosted
in the same country and tested with dig -
\url{http://linux.die.net/man/1/dig}}, the average resolution time is
acceptable within the metrics for international queries. A special case is
presented on the network tagged as \textit{``4Low''} were the time that take
to effectively resolving queries to any website is extremely high but after
the resolution finish,the network looks as had no major problem in data
acquisition. Also no mayor relation can be seen between the bandwidth and the
DNS resolve time.

\begin{table}[ht]
\begin{center}
\begin{tabular}{|c||c||c|c||c|c||}
 \hline
& & \multicolumn{2}{|c||}{Buffer} & \multicolumn{2}{|c||}{Performance} \\ \hline
Location	& DNS (ms) 	& Uplink (ms)	& Downlink (ms) & Latency (ms)	& Loss (\%) \\ \hline \hline
casa\_w		& 190		& 290			& 190 			& 140			& 0 		\\ \hline
casa2\_w	& 180		& 280			& 180 			& 140			& 0 		\\ \hline
casa\_c		& 180		& 280			& 190			& 140			& 0 		\\ \hline
marcia\_w	& 190		& 1200			& 860			& 190			& 0,5 		\\ \hline
marcia\_c	& 180		& 990			& 2100			& 190			& 0 		\\ \hline
sara\_w		& 220		& 5100			& 470			& 160			& 4 		\\ \hline
sara\_c		& 200		& 5100			& 459			& 160			& 4 		\\ \hline
polmos\_w	& 220		& 360			& 160			& 180			& 0 		\\ \hline
polmos2\_w	& 180		& 370			& 160			& 190			& 0 		\\ \hline
nalucem\_w	& 210		& 5100			& 1800			& 160			& 1,5 		\\ \hline
nalucem2\_w	& 210		& 5100			& 1800			& 210			& 0,2 		\\ \hline
mbahamon\_w	& 200		& 2900			& 290			& 200			& 1,5 		\\ \hline
4low\_w		& 1300		& 2900			& 590			& 180			& 0,5 		\\ \hline
hidalgo\_w	& 180		& 260			& 100			& 200			& 1,5  		\\ \hline
\end{tabular}
\caption[Netalyrz Test:DNS resolution time, Buffer time and Performance]{DNS resolution time, Buffer time and Performance}
\label{table:buffer}
\end{center}
\end{table}

The next four lines are crucial to proof existence of the phenomenon. The
buffer section can reveal how much time a packet spent in the existing buffers
along the way or, in other words into the link between the network and the
server. Unfortunately, in the uplink times are about $\sim300ms$, which
according by ICSI are times that in some cases may present a degraded
performance (as in online games or real time conference). For the downlink,
the networks that already are marked with high buffering time in the uplink
are the same for the opposite route. While the high measured times in the
uplink could it be justified by the diminish capacity to put new data in the
network against the capacity related with the downlink (related with the
asynchronous bandwidth capacity) there are cases in which the time is
excessive and Netalyzr alert the presence of excessive buffers possibly
generated by Bufferbloat.

\begin{table}[ht]
\begin{center}
\begin{tabular}{|c||c|c||}
 \hline
 & \multicolumn{2}{|c||}{Bandwidth (Mbps)} \\ \hline 
Location 	& Uplink & Downlink			   \\ \hline \hline
casa\_w		& 5		 & 14 				   \\ \hline
casa2\_w	& 5		 & 14 			   	   \\ \hline
casa\_c		& 5		 & 15 			       \\ \hline
marcia\_w	& 0,47	 & 1,5 			   	   \\ \hline
marcia\_c	& 0,54	 & 6,2 			   	   \\ \hline
sara\_w		& 0,57	 & 6,3 			   	   \\ \hline
sara\_c		& 0,57	 & 6,3 			   	   \\ \hline
polmos\_w	& 2,1	 & 9,6 			   	   \\ \hline
polmos2\_w	& 2,1	 & 9,6 			   	   \\ \hline
nalucem\_w	& 0,57	 & 4 			   	   \\ \hline
nalucem2\_w	& 0,57	 & 4 			   	   \\ \hline
mbahamon\_w	& 0,54	 & 7,8 			   	   \\ \hline
4low\_w		& 0,54	 & 6,6 			   	   \\ \hline
hidalgo\_w	& 1		 & 11 			   	   \\ \hline
\end{tabular}
\caption[Netalyrz Test: Bandwidth]{Bandwidths measured with Netalyzr}
\label{table:Bandwidth}
\end{center}
\end{table}


The main characteristic of Bufferbloat is the high latency, with some effect
on the packets loss, since packets spends most of the time in buffers along
the route. Netalyzr actually reveals that in at least two networks,
\textit{marcia} and \textit{sara}, this phenomenon occurs because although the
latency is low under normal conditions there is packet loss and buffer times
higher than normal. The packet loss can also be due to the conditions of the
experiments, and because due its characteristics, it is more common higher
packet lost using wireless but for these experiments there were no obstacles
between the router and computer located less than a meter from each other. Also
it is important to remember that the way of buffering time and latency are
calculated is through two different experiments with low and high load on the
network.

Among the characteristics of Netalyzr is the capacity to introduce a summary
table where it can be seen the different problems that arise in the networks
and the severity thereof. In Table \ref{table:Errores} are presented the
summary after all networks.

\begin{table}
\centering
\caption[Netalyrz Test: Summary of errors and warnings in Netalyzr]{Summary of errors and warnings in Netalyzr}
\label{table:Errores}
\subcaption*{Panel A: Summary or errors}
\begin{tabular}{|c||c|}
 \hline
Location		& Alerts            \\ \hline \hline
casa\_w			& 1,2,8             \\ \hline
casa2\_w		& 1,2,8             \\ \hline
casa\_c			& 1,2,8 			\\ \hline
marcia\_w		& 1, 3, 6, 9 		\\ \hline
marcia\_c		& 1, 3, 6, 8, 9 	\\ \hline
sara\_w			& 1, 2, 4, 5, 6, 8 	\\ \hline
sara\_c			& 1, 2, 4, 6, 8 	\\ \hline
polmos\_w		& 1,8 				\\ \hline
polmos2\_w		& 1, 3, 6, 8 		\\ \hline
nalucem\_w		& 1, 2, 6, 8 		\\ \hline
nalucem2\_w		& 1, 2, 6, 8, 10 	\\ \hline
mbahamon\_w		& 1, 3, 6, 8, 10 	\\ \hline
4low\_w			& A, 1, 3, 6, 7, 8 	\\ \hline
hidalgo\_w		& 1, 8 				\\ \hline
\end{tabular}
\bigskip
\subcaption*{Panel B: Description of errors}
\begin{tabular}{|c||c|}
 \hline
 \multicolumn{2}{|c|}{Alerts} \\ \hline \hline
A	& ISP's DNS is slow to lookup names								\\ \hline
1	& Certain TCP protocols are bloqued in outbound traffic 		\\ \hline
2	& The network does not reply when it needs fragmented traffic 	\\ \hline
3	& Fragmented UDP traffic if bloqued 							\\ \hline
4	& The packet loss was somewhat high 							\\ \hline
5	& The time to set up a TCP connection was somewhat high 		\\ \hline
6	& Network packet buffering is excessive 						\\ \hline
7	& DNS resolver may have probelms with DNSSEC 					\\ \hline
8	& Only some root servers returned DNSSEC information 			\\ \hline
9	& Not all DNS types were correctly processed 					\\ \hline
10	& The network indicated bursts of packet loss 					\\ \hline
\end{tabular}
\end{table}

With this summary is easy to clarify that the networks mentioned noted
actually have problems with the time with their buffers, specifically troubled
networks are: \textit{marcia, sara, nalucem, mbahamon} y \textit{4low}.
Coincidentally, less bandwidth networks are the same with longer buffers. In
addition, as mentioned, the network \textit{4low} also has serious issues
with the name resolution.


%\begin{table}[ht]
%\begin{center}
%\begin{tabular}{|c||c|}
% \hline
%Location		& Alerts            \\ \hline \hline
%casa\_w			& 1,2,8             \\ \hline
%casa2\_w		& 1,2,8             \\ \hline
%casa\_c			& 1,2,8 			\\ \hline
%marcia\_w		& 1, 3, 6, 9 		\\ \hline
%marcia\_c		& 1, 3, 6, 8, 9 	\\ \hline
%sara\_w			& 1, 2, 4, 5, 6, 8 	\\ \hline
%sara\_c			& 1, 2, 4, 6, 8 	\\ \hline
%polmos\_w		& 1,8 				\\ \hline
%polmos2\_w		& 1, 3, 6, 8 		\\ \hline
%nalucem\_w		& 1, 2, 6, 8 		\\ \hline
%nalucem2\_w		& 1, 2, 6, 8, 10 	\\ \hline
%mbahamon\_w		& 1, 3, 6, 8, 10 	\\ \hline
%4low\_w			& A, 1, 3, 6, 7, 8 	\\ \hline
%hidalgo\_w		& 1, 8 				\\ \hline
%\end{tabular}
%\caption[Netalyrz Test: Summary of errors and warnings in Netalyzr]{Summary of errors and warnings in Netalyzr}
%\label{table:Errores}
%\end{center}
%\end{table}

%\begin{table}[ht]
%\begin{center}
%\begin{tabular}{|c||c|}
% \hline
% \multicolumn{2}{|c|}{Alerts} \\ \hline \hline
%A	& ISP's DNS is slow to lookup names								\\ \hline
%1	& Certain TCP protocols are bloqued in outbound traffic 		\\ \hline
%2	& The network does not reply when it needs fragmented traffic 	\\ \hline
%3	& Fragmented UDP traffic if bloqued 							\\ \hline
%4	& The packet loss was somewhat high 							\\ \hline
%5	& The time to set up a TCP connection was somewhat high 		\\ \hline
%6	& Network packet buffering is excessive 						\\ \hline
%7	& DNS resolver may have probelms with DNSSEC 					\\ \hline
%8	& Only some root servers returned DNSSEC information 			\\ \hline
%9	& Not all DNS types were correctly processed 					\\ \hline
%10	& The network indicated bursts of packet loss 					\\ \hline
%\end{tabular}
%\caption[Netalyrz Test: Description of errors and warnings in Netalyzr]{Description of errors and warnings in %Netalyzr}
%\label{table:DescErrores}
%\end{center}
%\end{table}

\subsection{Iperf test}
Graphs and what happend with the rtts


\subsection{Benchmark test}
While ordinary Internet users request mainly optimized  web sites to make the
most of each user connection, it is natural that the traffic generated by
several requests of web sites simultaneously generates some kind of significant
load to the network. Or so it is what naturally should think.

Unfortunately with the results obtained with this Page Benchmark it is more
clear to see it is not so. By analyzing Figure \ref{fig:loadmeans}, the
networks that already were categorized as \textit{with problems}, follow that
trend and even with all the available bandwidth, takes longer to complete the
request. As example, the network \textit{marcia\_w} without other sources of
loads took a little over 9 seconds to complete the site request \footnote
{Again, the website is \url{http:// www.usm.cl}}. This results is way above
expectations since under normal conditions, the normal average is between 1-7
seconds with a tendency to be close to 2 seconds). An important study case of
study is the network \textit{mbahamon\_w} which was one of the lowest overall
in cases without overloading the network. While on a later date after these
tests, it was necessary to change the router used as access point, which
generated a lot of noise and interference for testing between tests seeing
themselves in a bad state at that time.

\begin{figure}[ht]
\centering
	%\rule{5.5cm}{7.1cm}
    \includegraphics[width=0.9\textwidth]{img/measures_page}
\caption[Page Benchmark: Total Load Means]{ Total Load means (in ms) based on \ref{app:pbmeasures}, Table \ref{table:loadmeans}}
\label{fig:loadmeans}
\end{figure}%

Before realizing that the router was the flaw in the network
\textit{mbahamon\_w}, under conditions of stress, this network was one of the
most affected networks delaying average time 23 seconds. The networks that are
in worst situations are those mentioned \textit{marcia\_w} and
\textit{sara\_c}. Due to the presence of large buffers which store packets of
uncontrolled manner and for longer periods, when the network's load decreases,
this resource  still under utilization. This is reflected in the observations
from from 3 onwards.

\begin{table}[ht]
\begin{center}
\begin{tabular}{|c||c|c|c||}
 \hline
Name 			& min		& max		& ratio  	\\ \hline \hline
4low\_w 		& 3415,1	& 20932,3	& 612,93374 \\ \hline 
casa\_w 		& 1654,8	& 3053,6	& 184,52985 \\ \hline 
casa2\_w		& 1646,6	& 2354,4	& 142,98555	\\ \hline 
casa\_c			& 1165,3	& 1944,1	& 166,83258	\\ \hline 
mbahamon\_w		& 2182,6	& 23868,5	& 1093,58105\\ \hline 
marcia\_w		& 9401,6	& 60119,0	& 639,45499	\\ \hline 
marcia\_c		& 10028,3	& 11140,0	& 111,08563	\\ \hline 
nalucem\_w 		& 6280,8	& 26048,6	& 414,73379	\\ \hline 
nalucem2\_w 	& 6349,4	& 32151,5	& 506,37068	\\ \hline 
polmos\_w 		& 1623,4	& 3701,4	& 228,00296	\\ \hline 
polmos2\_w 		& 1738,5	& 4861,8	& 279,65487	\\ \hline 
sara\_w 		& 2808,6	& 26561,5	& 945,72029	\\ \hline 
sara\_c 		& 2757,6	& 70423,3	& 2553,78953\\ \hline 
\end{tabular}
\caption[Page Benchmark: Minimum, Maximum and variation ratio.]{Minimum, Maximum and variation ratio.}
\label{table:varatio}
\end{center}
\end{table}

In Table \ref{table:varatio}, can be seen the comparison between the minimum
and maximum average load value across over all iterations on the same network
and the ratio between these values. While the ratio can not be compared
directly between two different networks, but can compare how many times the
maximum value was respect to its original value \footnote{a network may have a
low ratio with times that have a low variance and high average time}. Here,
the network \textit{marcia\_c}, and in the Table \ref{table:loadmeans},
accomplish times that were close to 11 seconds for all five iterations. With
this values, the ratio is low (11\%) but the real values are relatively high
for normal conditions (around 11 seconds).

Again, the networks \textit{marcia\_w}, both cases for \textit{sara}, stand
here because the high level of variation around 6 to 20 times the lowest
average time; highlighting \textit{mbahamon\_w} and \textit{sara\_c}. Contrary
to what was expected, in the case of the wired iteration over sara,
\textit{sara\_c}, the maximum variation was significantly higher in this
iterations than for those observed using wireless technology.

\begin{table}[!ht]
\begin{center}
\begin{tabular}{|c||c|c|c|c|c||}
 \hline
 & \multicolumn{5}{|c|}{ Ratio (\%)} \\ \hline
Name 		& 1			& 2			 & 3	        & 4				& 5 			\\ \hline \hline
4low\_w		& 155,55556	& 707,62342	 & 196,64754	& 232,87492		& 3302,32392	\\ \hline
casa\_w		& 108,35366	& 299,26429	 & 116,08775	& 356,49786		& 118,85296 	\\ \hline
casa2\_w	& 113,99878	& 281,65249	 & 127,18041	& 103,10219		& 117,71533 	\\ \hline
casa\_c		& 112,57589	& 304,73684	 & 103,63636	& 103,57766		& 104,58478 	\\ \hline
mbahamon\_w	& 163,95912	& 183,86356	 & 153,77847	& 139,55813		& 319,52128 	\\ \hline
marcia\_w	& 309,24808	& 6298,01469 & 267,57064	& 198,10685		& 205,48759 	\\ \hline
marcia\_c	& 310,36182	& 187,34044	 & 146,11202	& 205,08497		& 138,25129 	\\ \hline
nalucem\_w	& 261,95410	& 254,73473	 & 109,85803	& 106,96041		& 105,23376 	\\ \hline
nalucem2\_w	& 980,66263	& 315,95548	 & 254,12946	& 284,86443		& 110,70346 	\\ \hline
polmos\_w	& 106,36537	& 263,10680	 & 3902,47219	& 8443,23607	& 4150,06459 	\\ \hline
polmos2\_w	& 159,75460	& 1991,19249 & 8206,53009	& 417,45121		& 8226,29199 	\\ \hline
sara\_w		& 402,38095	& 240,75420	 & 394,93299	& 976,89312		& 395,52042 	\\ \hline
sara\_c		& 111,04952	& 281,07638	 & 109,03226	& 107,01366		& 109,54898 	\\ \hline
\end{tabular}
\caption[Page Benchmark: Ratio over own iteration]{Ratio over own iteration.}
\label{table:variationratio}
\end{center}
\end{table}

Also, for some cases adding another source was not decisive to obtain times
quite high. As can be seen in \ref{table:variationratio}, for
\textit{polmos\_w}, which in its latest iteration had a maximum of 127.343 ms
\footnote{The detailed values are available with CD} and a minimum of 1548ms.
Grounds for these outlayers measurements may be due to many factors ranging
from the conection setup to server problems.


\subsection{Smokeping Test}
It is already clearly established the presence of bufferbloat at least three
networks, resulting in a decrease in the performance and usability of the
networks, resulting in catastrophic times of answer for anyone who wants to
have a smooth sesion with any other service that requires capacity of the
network.

Trying to assimilate otherwise these results is that in Figure
\ref{fig:smokenat} can see the contrast of the two opposite poles.

\begin{figure}
\begin{subfigure}{\textwidth}
  \centering
	%\rule{5.5cm}{7.1cm}
    \includegraphics[width=0.8\textwidth]{img/smoke_nat_good}
\caption[Smokeping: Ping test to National servers with good performance]{Good performance Network}
\label{fig:smokenatgood}
\end{subfigure}%
\\
\begin{subfigure}{\textwidth}
\centering
	%\rule{5.5cm}{7.1cm}
    \includegraphics[width=0.8\textwidth]{img/smoke_nat_bad}
\caption[Smokeping: Ping test to National servers with bad performance]{Bad performance Network}
\label{fig:smokenatbad}
\end{subfigure}
\caption[Smokeping: Ping test to National servers]{Ping test to National servers}
\label{fig:smokenat}
\end{figure}

Figure \ref{fig:smokenatgood} shows the behavior of the network
\textit{casa\_w} which has faster response times with  no load around 50ms and
to superor  200ms with maximum use of the network. While can be appreciate
greater amount of \textit{``smoke''} (bar with same color as the line), this
is because by the time range that is driving (very low), tends generates this
small variation.

In contrast, in Figure \ref{fig:smokenatbad}, the amount of smoke is much
lower, but the variation in time is much higher, resulting in lower time to
100ms before loading, then vary drastically reaching over two seconds. While
this case was less packet loss, test took almost the double because of the
tendency to vary dramatically in the second case.


\begin{figure}
\begin{subfigure}{\textwidth}
\centering
	%\rule{5.5cm}{7.1cm}
    \includegraphics[width=0.8\textwidth]{img/smoke_int_good}
\caption[Smokeping: Ping test to International Servers with good performance]{Good performance network}
\label{fig:smokeintgood}
\end{subfigure}%
\\
\begin{subfigure}{\textwidth}
\centering
	%\rule{5.5cm}{7.1cm}
    \includegraphics[width=0.8\textwidth]{img/smoke_int_bad}
\caption[Smokeping: Ping test to International servers with bad performance]{Bad performance network}
\label{fig:smokeintbad}
\end{subfigure}
\caption[Smokeping: Ping test to International servers]{Ping test to International servers}
\label{fig:smokeint}
\end{figure}

For international servers, it can be appreciated that the tendency remains
very strong. The times related to Figure \ref{fig:smokeintbad} increases to 8
times its minimum value, ranging from 250ms to about 2 seconds. Again, the
loss in Figure \ref{fig:smokeintgood} is almost double than in Figure
\ref{fig:smokeintbad}, but the latter spent a longer period without load,
which reflect that in average have lower loss.

It is also interesting to note that while both servers are in the U.S., the
requests to the server noticed by the orange line took almost twice the time
against the blue-line-host  (above 300ms). This may be due to the
characteristics of the service provided or the resolution time by the server
because after a little more investigation, it was determined that was hosted
on a VPS with characteristics similars to Iperf's server connection used (same
hosting service).

\begin{figure}
\begin{subfigure}{\textwidth}
\centering
	%\rule{5.5cm}{7.1cm}
    \includegraphics[width=0.8\textwidth]{img/smoke_inf_good}
\caption[Smokeping: Web requests with good performance]{Good performance network}
\label{fig:smokewebgood}
\end{subfigure}%
\\
\begin{subfigure}{\textwidth}
\centering
	%\rule{5.5cm}{7.1cm}
    \includegraphics[width=0.8\textwidth]{img/smoke_inf_bad}
\caption[Smokeping: Web requests with bad performance]{Bad performance network}
\label{fig:smokewebbad}
\end{subfigure}
\caption[Smokeping: Web requests to National servers]{Web requests to National servers}
\label{fig:smokeweb}
\end{figure}

To conclude, Figure \ref{fig:smokeweb}  reveal the results of the HTTP
requests. In the case of Figure \ref{fig:smokewebgood}, load average achieved
was 1.1 seconds  with a loss of approximately 1.14\%. In contrast, in Figure
\ref{fig:smokewebbad}, right after the first overload (close at 17:00 hrs) of
the network, communication is affected, only been able to resume the
communication upon completion of the test. In the second iteration, becomes
not only intermittent communication, but also when it generate some traffic,
it takes about 3 seconds and has a loss of about 70\%. This not only affirms
our theory about the behavior and relationship to loss over time that lasted
the test, but also shows how terrible it would be for anyone trying to load
any website.


\subsection{Summarising Results}
The results are synthesized in order to analyze them comprehensively. 
The first test gave us the data to define the bandwidth of each
network with which subsequently went to work. Overall the upload link (uplink)
behaved as expected, performing in some cases gains of up to 50\% or more than
contracted. Obviously this could be due to generally low values ($\sim1Mbps$)
are therefore at a small variance is expressed as a high variation ratio. For
networks with greater uplink, between 1 and 5 Mbps also showed good
performance. On the other hand, the downlink reveal more variation between the
contracted and obtained. The circumstantial differences between a network with
good behavior to those with certain problems got reflected at this point. A
special case was the network \textit{polmos}. Strangelym after being tested with two
different Operating Systems (but on the same computer) results near
the contracted value were never obtained; not with other devices, where exhibit
values around the 40Mbps. The range values for the ping was not as expected,
but they still accepted values without being detrimental to smooth navigation
with a use that does not cause an excessive stress to the system.

With Netalyzr it is possible to obtain a clearer idea of the fundamentals and the
inner workings of networks. Also most revealing information was obtained from
the theoretical perspective, the existence of Bufferbloat obtaining
approximate times of the duration of the buffers present in these networks.
This results was the most important for the continuity of the study, mainly
because it proceeded to rule out or find problems associated with other
factors, for example a defective router (as in the case of \textit{mbahamon})
or environmental interference.

Iperf revealed to us the behavior of the networks trying to use their full
capacity. Also checked the behavior of the protocols and the different phases
that they have (slow start and congestion avoidance periods). While this was
not within the objectives of the study, witness the behavior helped to verify
that the operating level protocols, were working as expected. Along with this,
it helped to ensure that the presence of buffers is lethal when trying to make
a fluid exchange data with a high level of network utilization.

The simulation of the experience gained by a user when trying to load a
website under periods of low and high loads also can enable to understand
better the effects of Bufferbloat. While the trend was already well defined by
the above tests, individual results included in the study were obtained, but
defined as outlayers, such as occurred in the network \textit{polmos\_w} where
a maximal value was obtained in the last iteration, which were expected  to be
more alike to the first iteration.

If it is believed that due to Bufferbloat the times was increased, thanks to
smokeping got demonstrated that not only increase considerably, but almost
impossible to web surf or almost very difficult to think in online games with
good performance. The resulting loss and increased latency are considerable
reaching about 12 times the initial value, and growing to 100\% loss
communication.

%end section Test
\newpage
%section Test
\section{Conclusions and Further Work}
conclusions

http://lalith.in/2012/02/15/fun-with-tcp-cubic/
Examples of iperf and wireless 

\cite{MathisMacroCAA}

further

analysis in specific of how the three different flows (Elephant, mice and ant)\cite{HaElephants} \cite{evolvshortlongflows} reacts to the bufferbloat

How the linux kernel now reacts to bufferbloat. \cite{TokeLinux}

The effects under laying the use of CTCP\cite{Tan06compoundtcp}\cite{4146841} in Windows and the recommendations of the Bufferbloat community\cite{windowstips}

ver notas en \cite{MathisMacroCAA}

cerowrt project \cite{cerowrt}

Also from \cite{main:ref:1} and \cite{Vu-Brugier}, we can see that the rule-of-thumb doesn't longer apply to backbone routers, and a better estimator of the size of a buffer with n flows would be no more than $B = (\overline{RTT}xC)/\sqrt{n}$. With the assumption that short-flows plays a very small effect, and that the buffer size is dictated by the number of long flows, this factor will be proof that routers are much longer than they need to be, possible by two order of magnitude.
%end section Conclusions
\newpage
%%%%%%%%%%%%%%%%%%%%%%%%%%%%%%%%
\begin{appendix}
\section{Definitions}\label{app:definitions}
\begin{description}

\item [Bandwidth Delay Product (BDP)]:\hfill \\  

\end{description}

\section{Speed Test results}\label{app:bwmeasures}
\begin{table}[ht]
\caption[Speed Test: Speeds and Ping measured]{Speeds and Pings measured. }
\begin{center}
\begin{tabular}{|c||c|c|c|c|}
\hline
Location  & Medium  & Uplink (Mbps) & Downlink (Mbps) & Ping (ms)\\ \hline\hline
 casa\_w     & Wireless & 4,95           & 15,48         & 18 \\ \hline       % mov 5x15%
 casa2\_w    & Wireless & 5,09           & 13,46         & 20 \\ \hline       % mov 5x15%
 casa2\_c    & Cable    & 5,15           & 15,44         & 19 \\ \hline       % mov 5x15%
 marcia\_w	& Wireless & 0,83	        &  6,3			& 29 \\ \hline
 marcia\_c	& Cable    & 0,83	        &  6,29	        & 29 \\ \hline
 sara\_w	    & Wireless & 0,56	        &  5,45	        & 31 \\ \hline
 sara\_c	    & Cable    & 0,57	        &  8,08	        & 30 \\ \hline
 polmos\_w   & Wireless & 2,32           & 19,61         & 26 \\ \hline       % vtr 2x40%
 polmos2\_w  & Wireless & 2,34           & 18,72         & 22 \\ \hline       % vtr 2x40%
 nalucem\_w  & Wireless & 0,57           & 3,97          & 35 \\ \hline       % mov 0,5x4%
 nalucem2\_w & Wireless & 0,56           & 4,06          & 36 \\ \hline       % mov 0,5x4%
 mbahamon\_w & Wireless & 0,6            & 10,21         & 24 \\ \hline       % vtr 0,5x10%
 4low\_w     & Wireless & 0,6            & 8,03          & 29 \\ \hline       % vtr 0,5x10%
 hidalgo\_2  & Wireless & 1,19           & 17,96         & 19 \\ \hline       % vtr 1x20%
\end{tabular}
\end{center}
\label{table:speeds}
\end{table}

\begin{table}[ht]
\caption[Speed Test: Measured and Contracted Speed ]{Measured and Contracted Speed.}
\begin{center}
\begin{tabular}{|c|c||c|c||c|c||}
\hline
&&\multicolumn{2}{|c||}{Measured (Mbps)} & \multicolumn{2}{|c||}{Contracted (Mbps)} \\ \hline	
Location	& Provider	& Uplink 	  	 & Downlink  &	Uplink	& Downlink \\ \hline \hline 
casa\_w		& Movistar	& 4,95			 & 15,48	 & 5		& 15 \\ \hline 
casa2\_w	& Movistar	& 5,09			 & 13,46	 & 5		& 15 \\ \hline 
casa\_c	 	& Movistar	& 5,15			 & 15,44	 & 5		& 15 \\ \hline 
marcia\_w	& VTR		& 0,83			 & 6,3 	 	 & 0,5		& 10 \\ \hline 
marcia\_c	& VTR		& 0,83			 & 6,29	 	 & 0,5		& 10 \\ \hline 
sara\_w		& Movistar	& 0,56			 & 5,45		 & 0,5		& 10 \\ \hline 
sara\_c		& Movistar	& 0,57			 & 8,08		 & 0,5		& 10 \\ \hline 
polmos\_w	& VTR		& 2,32			 & 19,61	 & 2		& 40 \\ \hline 
polmos2\_w	& VTR		& 2,34			 & 18,72	 & 2		& 40 \\ \hline 
nalucem\_w	& Movistar	& 0,57			 & 3,97	 	 & 0,5		& 4 \\ \hline
nalucem2\_w	& Movistar	& 0,56			 & 4,06		 & 0,5		& 4 \\ \hline
mbahamon\_w	& VTR		& 0,6 			 & 10,21	 & 0,5		& 10 \\ \hline
4low\_w		& VTR		& 0,6			 & 8,03		 & 0,5		& 10 \\ \hline
hidalgo\_w	& VTR		& 1,19			 & 17,96	 & 1		& 20 \\ \hline
\end{tabular}
\end{center}
\label{table:comparative}
\end{table}

\newpage
\end{appendix}
%%%%%%%%%%%%%%%%%%%%%%%%%%%
%incluyo referencias
\bibliography{references,rfc}{}
%%%%%%%%%%%%%%%%%%%%%%%%%%%%

\end{document}
